\documentclass[11pt]{article}
\usepackage{fontspec}
\usepackage[french,english]{babel}
\usepackage{amsmath}
\usepackage{amssymb}
\usepackage{amsthm}
\usepackage{mathrsfs}
\usepackage{tikz-cd}
\usepackage{enumitem}

\newcommand{\printedpage}[1]{\marginpar{\footnotesize p.~#1}}
\newcommand{\msp}[1]{\marginpar{\footnotesize [#1]}}
\newcommand{\Ch}{\underline{\mathrm{Ch}}}
\newcommand{\Hom}{\underline{\mathrm{Hom}}}
\newcommand{\Fl}{\mathrm{Fl}}
\newcommand{\card}{\mathrm{card}}
\newcommand{\uHom}{\underline{\mathrm{Hom}}}

\begin{document}

\printedpage{9}
% head: A. Grothendieck : Dérivateurs, Ch. V, Édition des Universités de Montpellier II et Paris VII

\msp{page 13}
which on passing to ind-objects gives back the preceding morphism, or, what amounts to the same thing, such that for every $m,n\in\mathbf{N}$ (of the same parity), there exists an $N\geq m+n$ large enough that we have a diagram
\[
\begin{tikzcd}
\Ch_{\Sigma,m}(X;x,y)\times\Ch_{\Sigma,n}(X;y,z) \arrow[r,"*_{m,n}"] \arrow[dd,hook] \arrow[dr] & \Ch_{\Sigma,m+n}(X;x,z) \arrow[d,hook] \\
 & \Ch_{\Sigma,N}(X;x,z) \arrow[d,hook] \\
\Ch_{\Sigma}(X;x,y)\times\Ch_{\Sigma}(X;y,z) \arrow[r,"?"] & \Ch_{\Sigma}(X;x,z)
\end{tikzcd}
\]
and where the upper triangle is \textit{commutative in} Cat\,hot. (One must suppose that for $N$ large enough, the composite
\begin{tikzcd}[cramped, sep=small]
{} \arrow[d,hook] & \\
{} \arrow[r,"?"] & {}
\end{tikzcd}
factors through $\Ch_{\Sigma,N}(X;x,z)$ \dots)
To see that the composition functors $*_{m,n}$ define a homomorphism of the ind-objects in (Cat\,hot), it suffices to prove it in the particular cases $(m',n')=(m+1,n+1)$. We are going to draw the figure assuming $m$, $n$ \textit{odd.}

I realize here that I have stupidly botched it, the composition of two paths of lengths \textit{$m$, $n$ odd is}

\msp{page 14}
\textit{not a path} of length $m+n$, i.e. a diagram of type $\Phi_n$, but a diagram of \textit{dual} type $\Phi_n^o=\Psi_n$. For example, for $m=n=1$
\[
\underbrace{A_0\longleftarrow A_1\longrightarrow A_2}_{c}=\underbrace{B_0\longleftarrow B_1\longrightarrow B_2}_{c'}
\]
is of type $\Phi_2^o$! So as not to complicate life too much, we are going to restrict ourselves in the preceding developments to $m$, $n$ \textit{even,} hence paths of length at least 4. It is a little ridiculous! We are reduced to proving commutativity up to homotopy in the case where we leave $c'$ fixed and where we send $\Ch_{\Sigma,m}(X;x,y)$ into $\Ch_{\Sigma,m+2}(X;x,y)$, or conversely $c$ fixed and $\Ch_{\Sigma,n}(X;y,z)$ into $\Ch_{\Sigma,n+2}$ \texttt{[read:}$\Ch_{\Sigma,n+2}(X;y,z)$\texttt{]}. Let us take the first case, and write a path $c$ without negative indices (this is not a very happy notation!)
\[
A_0\longrightarrow A_1\longleftarrow A_2\cdots\longleftarrow A_{2m}
\]
(NB because of the condition $n$ even, all our paths now begin with an arrow from $A_0$ to $A_1$, never in the opposite direction).

I see that making the homotopy explicit, while offering no difficulty, is very tedious to explain. I am going to try to give a conceptual proof, without tedious ``calculation''.

% foot: Version 9 septembre 2001

\printedpage{10}
% head: A. Grothendieck : Dérivateurs, Ch. V, Édition des Universités de Montpellier II et Paris VII

\msp{page 15}
The two composite functors
\[
f=k\circ *_{m,n}\ ,\quad g=*_{m',n'}\circ(i\times j)
\]
in the square on page 11, which we have to prove are homotopic in $\uHom(\mathcal{A},\mathcal{B})$, can be viewed as obtained in the following way, in terms of two functors
\[
\varphi,\psi:\Phi_{m'+n'}\rightrightarrows\Phi_{m+n}, \tag*{(1)}
\]
whence two ``transposed'' functors
\[
\begin{tikzcd}
\uHom(\Phi_{m,n},X) \arrow[r,shift left,"F"] \arrow[r,shift right,"G"'] & \uHom(\Phi_{m',n'},X)
\end{tikzcd} \tag*{(2)}
\]
\texttt{[rather} $\Phi_{m+n}$\texttt{,} $\Phi_{m'+n'}$\texttt{]} which functors ``induce'' $f$ and $g$. To be precise, denoting by $a$, $b$ the source and target objects of $\Phi_{m+n}$, by $c$ the object which corresponds to the target of the path $c$ of length $m$ \texttt{[too many} $c$\texttt{'s]}, to the source of the path $c'$ of length $n$ which is to be composed, finally by $a_0$, $b_0$ the source and target elements of $\Phi_{m'+n'}$, the functors under consideration (1) send $a_0$ to $a$, $b_0$ to $b$, hence define a commutative diagram
\[
\begin{tikzcd}
\Phi_{m'+n'} \arrow[r,shift left,"\varphi"] \arrow[r,shift right,"\psi"'] & \Phi_{m+n} \\
\underset{\text{\footnotesize discrete category}}{\varepsilon=\{a_0,b_0\}.} \arrow[u,hook,"\alpha'"] \arrow[ur,hook,"\alpha"] &
\end{tikzcd} \tag*{(3)}
\]
This said, the target category $\mathcal{B}$ for $f$, $g$, namely $\Ch_{\Sigma,m'+n'}(X;x,z)$, is interpreted as the full subcategory of $\uHom(\Phi_{m+n},X)$ consisting

\msp{page 16}
of the functors $c':\Phi_{m'+n'}\longrightarrow X$ whose composite with $\alpha'$ is the functor
\[
\beta:\varepsilon\longrightarrow X\qquad\qquad
\begin{aligned}
a_0&\longrightarrow x\\
b_0&\longrightarrow z,
\end{aligned} \tag*{$(*)$}
\]
and likewise the source category $\mathcal{A}$ is interpreted as the full subcategory of $\uHom(\Phi_{m+n},X)$ consisting of the $c:\Phi_{m+n}\longrightarrow X$ whose restriction to $\{a,b,c\}$ is the functor
\[
a\longrightarrow x,\ b\longrightarrow y,\ c\longrightarrow z.
\]
\texttt{[rather} $b\longrightarrow z$\texttt{,} $c\longrightarrow y$\texttt{]} It is contained in the full subcategory $\mathcal{A}'$ of functors whose restriction to $\{a,b\}$ is $a\longrightarrow x,\ b\longrightarrow z$, i.e. whose composite with $\alpha$ is the functor $\beta:\varepsilon\longrightarrow X$ already considered above. Thus the functors $f,g:\mathcal{A}\rightrightarrows\mathcal{B}$ factor through $\mathcal{A}'$
% foot: Version 9 septembre 2001
\printedpage{11}
% head: A. Grothendieck : Dérivateurs, Ch. V, Édition des Universités de Montpellier II et Paris VII
\[
\begin{tikzcd}
\mathcal{A} \arrow[r,hook] \arrow[dr,shift left,"f"] \arrow[dr,shift right,"g"'] & \mathcal{A}' \arrow[d,shift right,"g'"'] \arrow[d,shift left,"f'"] \\
 & \mathcal{B},
\end{tikzcd}
\]
functors deduced from $\varphi$, $\psi$ above, giving rise to the commutative triangles (3), and from the functor $\beta:\varepsilon\longrightarrow X$, in the said way. To prove that $f$ and $g$ are homotopic, it suffices to prove it for $f'$, $g'$. And for this, it suffices to show that $\varphi$ and $\psi$ are homotopic as objects of the category $\uHom(\Phi_{m+n},\Phi_{m'+n'};\alpha)$ consisting of the functors whose restriction to $\varepsilon$ is $\alpha$ \texttt{[rather of the category} $\uHom(\Phi_{m'+n'},\Phi_{m+n};\alpha')$ \texttt{consisting of the functors whose restriction to} $\varepsilon$ \texttt{is} $\alpha'$ \texttt{]}.

\msp{page 17}
Moreover, the two assertions (on $f'$, $g'$, for arbitrary $X$, $x$, $y$, $z$, \texttt{[rather} $x$\texttt{,} $z$\texttt{]} and on $\varphi$, $\psi$, $\alpha$ \texttt{[rather} $\alpha'$\texttt{]}) are in fact equivalent as one sees by taking $X=\Phi_{m+n}$, $x=a$, $y=c$, $z=b$ \texttt{[rather} $x=a$\texttt{,} $z=b$\texttt{]}. This will be contained in the more general lemma. (\footnote[4]{In the considerations that follow, we in fact tacitly take the functors under consideration to be \textit{increasing} on the vertices with respect to the order of their numbering. Thus one should write $\uHom_{cr}$ instead of $\uHom$, in the statements that follow, including in the connectedness lemma opposite.})

\textbf{Lemma.} \textit{Let $k$ and $l$ be two integers. Then the full subcategory of $\uHom(\Phi_k,\Phi_l)$ consisting of the functors $\Phi_k\longrightarrow\Phi_l$ sending source to source, target to target, is connected.} (NB It is nonempty if and only if $k\geq l$.)

We are going to describe these functors completely (= increasing maps) $\Phi_k\xrightarrow{\varphi}\Phi_l$, identifying $\Fl(\Phi_k)$ \texttt{[rather the set of nonidentity arrows of} $\Phi_k$\texttt{]} with the interval of integers $[1,2k]$, and associating to $\varphi$ the subset
\[
E_\varphi\subset\Fl(\Phi_k)=[1,2k]\qquad\qquad E_\varphi=\{u\in\Fl(\Phi_k)\mid\varphi(u)\text{ an identity}\}
\]
consisting of the $u\in\Fl(\Phi_k)$ which are ``contracted'' by $\varphi$ to an identity. If $E$ is a set of integers, we introduce on it the notion of connected components (by the equivalence relation in $E$ generated by the successor relation $y=x+1$). One sees at once that the connected components of $E_\varphi$ which contain neither the first nor the last arrow of $\Phi_k$ are even, and that % [sic] connnexes -> connexes, triple n
\[
\underbrace{\card\Fl(\Phi_k)}_{2k}=\card\,(E_\varphi)+\underbrace{\card\Fl(\Phi_l)}_{2l},
\]

\msp{page 18}
i.e.
\[
\card\,(E_\varphi)=2(k-l).
\]

% foot: Version 9 septembre 2001

\printedpage{12}
% head: A. Grothendieck : Dérivateurs, Ch. V, Édition des Universités de Montpellier II et Paris VII

It is immediate that a functor $\varphi$ is surjective on objects and on arrows, hence identifies $\Phi_l$ with a quotient category of $\Phi_k$ - and the equivalence relation in $\Phi_k$ defined by $\varphi$ is the one obtained by identifying all the objects that are source or target of arrows of one and the same connected component $I$ of $E_\varphi$
\[
\cdots\longleftarrow\ \underbrace{\longrightarrow\ \longleftarrow\ \longrightarrow\ \longleftarrow}_{I}\ \longrightarrow\cdots
\]
and the arrows in this component with the corresponding identity. The conditions stated on $E$ (independently of a $\varphi$ which would give rise to it ensure that the quotient category is isomorphic (with \textit{preservation} of source and target) either to $\Phi_l$, or to the dual $\Psi_l=\Phi_l^o$, according to the direction of the first arrow of $\bar{\Phi}_k$ (quotient of $\Phi_k$). This ``initial direction'' is the same as that for $\Phi_k$ if the first arrow of $\Phi_k$ does not belong to $E$, or else if it belongs to it and its connected component in $E$ has even cardinality. In the opposite case ($u\in E$ and its connected component in $E$ has odd cardinality), the initial direction of $\bar{\Phi}_k$ is opposite to that of $\Phi_k$. % [sic] "(indépendamment d'un $\varphi$ qui lui donnerait naissance" -> parenthèse fermante manquante, the opening parenthesis is never closed; [sic] connnexe -> connexe, triple n

\msp{page 19}
One sees therefore that $E$ must satisfy yet another condition, for the corresponding $\bar{\Phi}_k$ to be isomorphic to $\Phi_l$, i.e. for $E$ to come from a functor $\varphi$. I summarize all the conditions

\begin{itemize}
\item[a)] $\card E=2(k-l)$.
\item[b)] The connected components of $E$ which contain no initial or final arrow have even cardinality.
\item[c)] If $\Phi_k$ and $\Phi_l$ have the same initial direction ($k$ and $l$ of the same parity), then the possible component of $E$ containing the initial arrow must itself also have even cardinality (which, because of a) \texttt{[rather a) and b)]}, is equivalent to the same condition for the possible connected component containing the final arrow of $\Phi_k$). If $\Phi_k$ and $\Phi_l$ have different initial directions, the initial arrow of $\Phi_k$ must belong to a connected component of \textit{odd} cardinality in $E$ (which is equivalent to likewise for the final arrow).
\end{itemize}

On the other hand, the functor $\varphi$ is uniquely determined by the knowledge of $E$, for there is a \textit{unique} isomorphism $\bar{\Phi}_k\xrightarrow{\sim}\Phi_l$ compatible with sources and targets, because of the

\textbf{Lemma} (immediate) : \textit{An automorphism of $\Phi_l$ which is the identity on the source object (or on the target object) is the identity everywhere.}

The functors $\varphi$, $\psi$ from earlier corresponded to the choice, one of $E_\varphi=\{1,2,k-1,k\}$ \texttt{[rather} $E_\varphi=\{1,2,2k-1,2k\}$\texttt{]}, the other of $E_\psi=\{1,2,i,i+1\}$, 2 connected components

\msp{page 20}
in one case, three \texttt{[rather two]} in the other.) % [sic] "l'autre.)" -> parenthèse fermante sans parenthèse ouvrante, no matching opening parenthesis

% foot: Version 9 septembre 2001

\printedpage{13}
% head: A. Grothendieck : Dérivateurs, Ch. V, Édition des Universités de Montpellier II et Paris VII

In the general case, one should make explicit the order relation on the set $\uHom^{!}(\Phi_k,\Phi_l)$ (! means : compatible with source and target), in terms of the representative sets $E_\varphi$, $E_\psi$. I am not going to make this order relation explicit, but only a property of it which will suffice for us. Let $\varphi_E$ be the functor defined by a subset $E\subset[1,2k]$. First of all :

\textbf{Lemma 1} (immediate) : \textit{Let $E,E'\subset[1,2k]$ be admissible, and let $j\in[1,2k]$ be such that $E\cap[1,j]=E'\cap[1,j]$. Then $\varphi_E$ and $\varphi_{E'}$ agree on the target of $u_j$ ($j$ an arrow of $\Phi_k$) and on all the objects of $\Phi_k$ preceding it (in the sense of the numbering of the objects, and not of the order relation of the category $\Phi_k$ of course.) Dually, if $E\cap[j,2k]=E'\cap[j,2k]$, then $\varphi_E$ and $\varphi_{E'}$ agree on the source of $y$, and on all the objects of $\Phi_k$ following it.} % [sic] la source de $y$ -> la source de $u_j$, apparently intended by duality with "le but de $u_j$"

Let us now consider a connected subset $I$ of $E$ (not necessarily a connected component), $I=[r,s]\subset[1,2k]$ ($r$ the smallest element of $I$, $s$ its largest element) and suppose that

\begin{itemize}
\item[a)] $\card I$ is even (i.e. $s-r$ is even) \texttt{[rather} $s-r$ \texttt{odd]}
\item[b)] $r\neq1$, and $r-1\notin E$ (i.e. $r-1\in[1,2k]\backslash E$)
\end{itemize}

and let
\[
E'=(E\backslash\{s\})\cup\{r-1\}=E\backslash I\cup\overbrace{[r-1,s-1]}^{I'}
\]
obtained from $I$ by ``translating'' one notch to the left.

\msp{page 21}
It is immediate that if $E$ is admissible, $E'$ is too (and conversely). This said, we have the

\textbf{Lemma 2.} \textit{$E$ and $E'$ being as above, there exists a homomorphism of functors between $\varphi_E$ and $\varphi_{E'}$} (I am not saying in which direction the arrow goes!). % [sic] une homomorphisme -> un homomorphisme, wrong gender

\textbf{Corollary.} Dual statement, when $E'$ is obtained from $E$ by translating to the right by one notch a connected interval of even length $I=[r,s]$ of $E$, $s\neq2k$ and $s+1\notin E$.

Indeed, in this case $E$ is obtained from $E'$ by an operation of the type made explicit in lemma 2.

The connectedness of the category $\uHom^{!}(\Phi_k,\Phi_l)$ is now clear. If $k$ and $l$ have the same parity, by pushing the connected components of $E$ to the left, we write $E$ in the ``normal'' form
\[
E_0=[1,2(k-l)].
\]
If $k$ and $l$ are of opposite parity, we cannot push the whole connected component of $E$ which contains the last element $2k$, for it has odd cardinality. One must leave $2k$ in $E$ and push the rest, and we reduce $E$ to the normal form
\[
E_0=[1,2(k-l)-1]\cup\{2k\}.
\]

\hfill q.e.d.!

% foot: Version 9 septembre 2001

\end{document}
